\section{Discussion}
\label{sec:discussion}

\subsection{Vertical placement of the canopy drag}

This study asked which modelled processes control transport in a shallow
vegetated lagoon, and on the fixed bed the answer the ensemble gives is the
seagrass canopy, with an effect five times that of wave coupling and the same
under both wave treatments. The mechanism the model offers for that is vertical
rather than frictional, since introducing the canopy removes momentum from the
lower part of the water column without removing a proportionate amount from the
upper part, more than doubling the ratio of surface to bed speed. Such a profile
is the signature of a canopy shear layer rather than of a rougher bed, which
acts through a logarithmic profile whose surface-to-bed ratio is bounded well
below the value obtained here \citep{Nepf2012}, and the bed shear stress
confirms the distinction by its sign, falling under the canopy where an
equivalent bed roughness raises it over the same cells. Since a surface drifter
samples the top of that profile, a treatment that reshapes the profile relocates
the drifter even where it changes the depth-averaged transport little.

That places the present result in a particular relation to the two earlier
models of this basin rather than in opposition to them.
\citet{Ingrassia2024} obtained a best-fit uniform Gauckler--Strickler
coefficient for the vegetated lagoon within a depth-averaged framework, where an
equivalent bed roughness is the appropriate representation of a canopy because
there is no vertical profile for the canopy to shape. \citet{DeMarchis2012}
worked in three dimensions and assigned Nikuradse roughness heights by seagrass
class, establishing the vertical structure of the wind- and tide-driven
circulation that motivates using a three-dimensional solver here at all. Both
answered the questions they set, and a roughness representation reproduces the
drag the meadow exerts. What has changed since is the availability of a
canopy formulation that resolves that drag over the vertical extent the plants
occupy, and the contribution of the present work is to quantify what that
capability buys in a basin of this depth. The canopy parameters used here were
taken from published measurement of the two species and were not adjusted
against any observation, so the improvement follows from where the drag is
applied rather than from how it was tuned.

\subsection{Information carried by each gauge position}

Water level separated the configurations by the canopy at AltaVilaEst, by wave
coupling at BocaNord and by neither at BocaSud, and the three positions
constrain different parts of the problem. Agreement at the inlet gauges
establishes that the offshore forcing, the datum and the exchange through the
two mouths are right, and in a micro-tidal basin with two open connections those
set most of the level there, the wave contribution to the northern inlet being
the one dynamic effect that reaches it. Agreement in the interior tests
something the outer pair cannot, because the level there is set by what the
water does on its way in, which is where the canopy acts. The roles are
complementary, and a match in the interior would be difficult to interpret
without the exchange having been established first.

Together with the drifter result, the gauge result refines the familiar argument
that water level constrains transport weakly and that Lagrangian observations
form a distinct constraint rather than a confirmation \citep{Liu2011,
Revelard2021}. Water level is not inherently insensitive to the canopy, since
the one gauge sited inside the vegetated basin identified the same factor as the
drifters and ranked the fixed-bed configurations the same way. For a basin of
this type the useful question is therefore less whether to validate on water
level than where to place the instrument intended to discriminate, and the
interior of the vegetated compartment is a considerably more informative
position for that purpose than the inlet. Section~\ref{sec:disc_synergy} shows
that the interior gauge also discriminates where the drifters do not.

\subsection{Bare-bed skill through a reshaped bathymetry}
\label{sec:disc_synergy}

The mobile-bed arm contains one configuration that, taken on its own, supports
a strong statement about process coupling. In the bare arm, wave coupling and
bed mobility each leave Lagrangian skill nearly unchanged or lower it when
applied alone and raise it when applied together, with an interaction term of
$+0.116$, tight bootstrap intervals and a mechanism available in the literature,
since wave orbital motion sets the bed stress that mobilises sediment and
wave--tide coupling is known to govern the sediment dynamics of shallow
micro-tidal basins \citep{Carniello2011}. The combination brings the bare
member to 0.708, within 0.03 of the vegetated members, and removes the canopy
contrast in that cell.

The segments restarted on 7~July show that the statement has no physical basis
in this model. The skill is carried by the bathymetry that six days of
wave-driven erosion reshaped, through a lagoon holding near
9~\si{\gram\per\litre} of suspended sediment and losing 4~million~\si{\metre\cubed}
of bed in nine days, a state the basin does not reach in summer. The reshaped
bed does not reproduce the vertical structure of the canopy either, the
surface-to-bed speed ratio of that member remaining at 2.17 against 5.13 with
the canopy, so the agreement in skill is reached through a different flow. A
model without a canopy can therefore reach the drifter skill of a model with
one through a mechanism that is wrong, and the drifters alone do not separate
the two. The interior gauge does, since at AltaVilaEst the bare member keeps an
anomaly RMSE of 0.040~m and a correlation of 0.86 against 0.024~m and 0.94 for
its vegetated counterpart.

The initial condition shows the same asymmetry by another route. A start from a
uniform offshore column lowers the skill of the bare members by 0.08 to 0.15 and
that of the vegetated members by at most 0.02, so the canopy contrast grows from
about 0.12 to about 0.19 when the start is wrong. In both cases an error
elsewhere in the model degrades the bare configurations and leaves the
vegetated ones nearly unaffected, and the magnitude of a canopy contrast
therefore depends on what else the model gets wrong, even where its sign does
not.

What follows from this bears on how such an ensemble is read rather than on this
basin alone, since canopy drag is commonly left out of lagoon and estuary models
or folded into a bed roughness coefficient, and an ensemble built that way can
still return a well-defined and statistically clean interaction between the
processes it does carry. Here that interaction is produced by an implausible
morphodynamic state, and nothing internal to the bare half of the design
distinguishes it from a physical result. What distinguished them was varying
the canopy as a factor in its own right, checking the morphodynamic state
against a plausible sediment budget, and scoring the result against two
observation types that no model decision had used.

\subsection{Limitations}
\label{sec:disc_limits}

\textbf{Wave dissipation by the canopy.} The canopy does not reach the wave
model, because the coupling passes a bed roughness field to the spectral model
while the canopy enters the momentum equation of the flow, so mean orbital
velocity at the bed changes by under 1\% when the canopy is introduced and not
in an order that follows canopy density. Seagrass meadows do attenuate waves, in
this species and at this scale \citep{Infantes2012, Prinos2010}, and shallow
restricted Mediterranean lagoons carry wave fields that respond to bed
conditions \citep{Morote2025}. The wave field in these simulations is therefore
that of a bare bed, the wave contrasts reported here are conditional on that,
and so is the wave-driven erosion that the canopy fails to shield in the
mobile-bed arm. Representing canopy dissipation inside the spectral model is the
corresponding extension. It would be expected to change the wave contrasts
rather than the canopy ones, which are measured on the flow.

\textbf{Canopy reconfiguration.} The canopy bends under load rather than holding
a fixed height, and the response is dynamic, the velocity ratio between two
probes differing only in leaf stiffness moving from 1.005 in the slowest
quintile of the record to 0.807 in the fastest 5\%, where a fixed shortening
would give a constant ratio. A 13\% change in leaf density moves the velocity
field by 20\% at the upper end of the range, which is a strong sensitivity to a
value taken from a published table rather than measured in this meadow. The
bending is moreover mechanical, whereas the flume measurements of
\citet{CiraoloFerreri2006} describe reconfiguration empirically through a
friction factor that varies by a factor of twelve over the velocity range they
tested. That empirical formulation is available in the model only in its
depth-averaged solver, so the two could not be compared directly and the
agreement between them was not tested.

\textbf{Bed mobility.} The morphodynamic arm is not calibrated, the critical
shear stress for erosion having been constrained by neither a repeated
bathymetric survey nor a turbidity record, and with waves it leaves any
plausible state of the basin, with suspended sediment near
9~\si{\gram\per\litre} and the lagoon bed lowered by 12~cm on average in nine
days (Section~\ref{sec:res_sensitivity}). It is therefore a sensitivity test and
is kept out of the attribution. Before it could enter it, the erosion
parameters would need to bring suspended sediment in the lagoon to the order of
tens of milligrams per litre.

\textbf{Window and restart segments.} The record covers two days of water level
and thirty-four hours of drifter observations within one summer season, and the
vegetated configurations are three-day restart segments. A segment run without
any other change reproduces its continuous member in water level and drifter
skill (Section~\ref{sec:res_wl}), but seasonal behaviour, storm conditions and
the interaction of a canopy with a mobile bed over longer periods are outside
what this design addresses.

\textbf{Initial condition.} The ensemble starts from the depth profiles of the
CMEMS field at the start time, which a start from a uniform column does not
approach within the nine days, and from the uniform start the canopy contrasts
on the fixed bed are $+0.195$ and $+0.188$ rather than $+0.120$ and $+0.114$
(Section~\ref{sec:res_sensitivity}). The stratified start is itself a single
profile applied over the domain and leaves out the horizontal structure of the
CMEMS field, so the canopy contrasts reported here may still carry part of a
start-up adjustment. The solution also depends on the host and the
decomposition at the Courant limit of the ensemble, by 0.034 in skill from the
uniform start, which all members share, while the horizontal and vertical
resolution do not limit the reported quantities.

\textbf{\emph{Cymodocea nodosa} parameters.} The \emph{C.\ nodosa} parameters
are loosely constrained, coming from a flume study of plants collected
elsewhere, since no measurement of that species in this lagoon was available.
The species carries about 1.5\% of the canopy frontal area in the model, which
bounds the consequence, and the drag is carried almost entirely by the
\emph{P.\ oceanica} column.
